\section{Conclusion and Limitations}
\label{sec:conclusion}

We presented \method{}: harness synthesis as a learned skill.  A LoRA
proposer, trained with a discovery-adapted group objective for 27 steps,
designs task-conditioned harnesses---prompts, skills, control parameters,
and reviewed executable tools---that carry a permanently frozen 9B executor
past a discovery-finetuned model of the same scale on all eleven tasks of a
heterogeneous suite, and match or exceed the published $\le$10B state of the art
on seven---setting a new best on six of them, including the Erd\H{o}s
minimum-overlap constant, where the frozen executor's own optimizer also passes
the best known human construction.  The campaign's instrumentation supports a reading beyond the
leaderboard: improvement decomposes into what the proposer learned (protocol
designs like the Hadamard parameter sweep), what the search found
(cascade-discovered program families, compounded by inheritance), and what
information unlocked (a timeout that favored no-ops; a scoring rule the
executor had never been told).  The last category carries the most transferable lesson: an apparent capability
ceiling of a small executor can be a legibility failure of its harness.  We note
that both such unblocks came from outside the learned loop---telemetry that we
read, and a curated note---so automating that diagnosis inside the proposer is a
target, not a result we demonstrate.

\paragraph{Limitations.}  Campaign results are single-seed (winner
stability was re-verified $5\times$; full campaigns were not replicated), and
two table entries carry provenance caveats stated in \S\ref{sec:exp-results}.
The comparison behind our headline result bundles harness search with proposer
training, so their relative contribution is unresolved until the
\method~(context) ablation lands.  Spontaneous tool authorship remains rare
without worked examples or forcing; a small SFT phase on tool-authoring
trajectories is the natural next lever.
Transfer of a trained proposer to unseen task families, and of its harnesses to
larger frozen executors---the amortization claim's strongest test---is under way
(\S\ref{sec:exp-transfer}, \S\ref{sec:exp-crossmodel}); scaling the proposer
beyond 9B is left to future work.

\paragraph{Safety.}  A system whose proposer writes code that runs inside
an evaluation loop must treat evaluation integrity as a design invariant,
not a hope.  \method{}'s defenses are structural: generated tools cannot read the
evaluator implementation or ground truth, repairs are held to capability parity
with the executor, and every drop, repair, and rejection is logged for
audit.  We release the review chain alongside the trainer, and view
leakage-controlled self-extension---not unrestricted self-modification---as
the appropriate operating point for automated agent design
today~\citep{dgm,skalse2022}.
